% 原书第六篇6.4节197--213（本节最后17题），印刷页255--256，
% PDF264--265；答案印刷348，PDF357。
\moleculeproblem{197}{255}
用零温自由电子态密度
$dn\propto\sqrt E\,dE$ 证明：
(a) 金属电子浓度为 $n$ 时最高动能 $E_F$；
(b) 零温平均电子动能与最高动能的比值。
\exsolution
三维含自旋的自由费米电子态密度为
\[
 dn=\frac{(2m_e)^{3/2}}{2\pi^2\hbar^3}
 \sqrt E\,dE,\qquad 0<E<E_F.
\]
积分得到
$n=(2m_eE_F)^{3/2}/(3\pi^2\hbar^3)$，
解得
\[
 \boxed{E_F=\frac{\hbar^2}{2m_e}(3\pi^2n)^{2/3}}.
\]
平均动能
\[
 \overline E=\frac{\int_0^{E_F}E\sqrt E\,dE}
 {\int_0^{E_F}\sqrt E\,dE}
 =\frac{(2/5)E_F^{5/2}}{(2/3)E_F^{3/2}}
 =\boxed{\frac35E_F}.
\]
\excheck 含自旋双重简并已计入态密度；
不能将费米电子平均动能误写成经典的 $3k_BT/2$。

\moleculeproblem{198}{255}
在 $T=0$ 的金属自由电子气中，
动能超过最大动能一半的电子占百分之多少？
\exsolution
零温仅填满 $0<E<E_F$ 的电子态，
其中 $dn\propto E^{1/2}\,dE$。
故
\[
 f=\frac{\int_{E_F/2}^{E_F}E^{1/2}dE}
 {\int_0^{E_F}E^{1/2}dE}
 =1-\left(\frac12\right)^{3/2}.
\]
所以
\[
 \boxed{f=1-2^{-3/2}\simeq0.6464
 \simeq64.6\%\approx65\%}.
\]
\excheck 这里积分的是\emph{态数}，
不是能量权重 $E\,dn$。

\moleculeproblem{199}{255}
钠在 $T=0$ 时费米能
$E_F=3.07\,\mathrm{eV}$，
密度 $\rho=0.97\,\mathrm{g/cm^3}$。
求平均每个钠原子贡献多少自由电子。
\exsolution
从
$E_F=\hbar^2(3\pi^2n_e)^{2/3}/(2m_e)$ 反解，
\[
 n_e=\frac1{3\pi^2}
 \left(\frac{2m_eE_F}{\hbar^2}\right)^{3/2}.
\]
钠原子数密度
$n_{\rm Na}=\rho N_A/M_{\rm Na}$，
其中 $M_{\rm Na}=22.99\,\mathrm{g/mol}$。
因此
\[
 \boxed{\nu=\frac{n_e}{n_{\rm Na}}
 =\frac{M_{\rm Na}}{3\pi^2\rho N_A}
 \left(\frac{2m_eE_F}{\hbar^2}\right)^{3/2}
 \simeq0.96\approx1}.
\]
\excheck 在模型精度下约为每原子一个自由电子；
书末简答约0.93，可能源于所采用的金属参数取值差异。

\moleculeproblem{200}{255}
把经典电子气加热到什么温度，
它的平均动能会与 $T=0$ 的铜金属自由电子的
平均动能相等？假设每个铜原子贡献一个自由电子。
\exsolution
自由电子数密度
$n=\rho_{\rm Cu}N_A/M_{\rm Cu}$。
零温平均动能 $\overline E_0=3E_F/5$，
其中
$E_F=\hbar^2(3\pi^2n)^{2/3}/(2m_e)$。
经典理想气体电子平均动能为
$3k_BT/2$，令两者相等，
\[
 \frac32k_BT=\frac35E_F
 \ \Longrightarrow\
 \boxed{T=\frac{2E_F}{5k_B}}.
\]
采用铜密度 $8.96\,\mathrm{g/cm^3}$
得到 $E_F\simeq7.04\,\mathrm{eV}$，
故
$\boxed{T\simeq3.27\times10^4\,\mathrm K
 \approx3\times10^4\,\mathrm K}$。
\excheck 该温度只是两个平均动能的\emph{数值等效温度}，
并非说铜在此温度仍保持固态金属电子气。

\moleculeproblem{201}{255}
$T=0$ 金属中自由电子浓度
$n=2.0\times10^{22}\,\mathrm{cm^{-3}}$，
金属体积 $V=1.0\,\mathrm{cm^3}$。
估算费米能附近相邻自由电子能级间隔（eV）。
\exsolution
在连续近似中粒子总数
$N(E)=V(2m_eE)^{3/2}/(3\pi^2\hbar^3)$。
费米能处态密度
$dN/dE=(3/2)(nV/E_F)$，
其中
$E_F=\hbar^2(3\pi^2n)^{2/3}/(2m_e)$。
因此平均相邻\emph{电子态}的能距
\[
 \boxed{\delta E\simeq
 \left(\frac{dN}{dE}\right)^{-1}_{E_F}
 =\frac{2E_F}{3nV}
 =\frac{\hbar^2}{3m_eV}(3\pi^2n)^{-1/3}}.
\]
数值按给定数据约为
$\boxed{\delta E\sim10^{-22}\,\mathrm{eV}}$。
\excheck 此定义计入每个空间轨道的两个自旋态；
若讨论\emph{不同能量值}而非电子态，
简并与边界几何会改变系数。原书给约
$2\times10^{-22}\,\mathrm{eV}$，与该量级一致。

\moleculeproblem{202}{255}
依据零温自由电子态密度，
(a) 推导电子速率分布 $dn_v$；
(b) 求平均速率与最大费米速率之比。
\exsolution
令 $E=m_ev^2/2$，
$0\le v\le v_F$，
由 $dn\propto\sqrt E\,dE$ 与
$dE=m_ev\,dv$，故
\[
 \boxed{dn_v=\frac{m_e^3}{\pi^2\hbar^3}v^2\,dv},
 \qquad
 \boxed{f(v)=\frac3{v_F^3}v^2},
 \quad 0<v<v_F.
\]
平均速率为
\[
 \overline v=\int_0^{v_F}vf(v)dv
 =\frac{3}{v_F^3}\frac{v_F^4}{4},
 \quad \boxed{\frac{\overline v}{v_F}=\frac34}.
\]
\excheck 零温下速率并非全部为零；
费米简并使电子填充从0到费米动量的态。

\moleculeproblem{203}{255}
从三维自由电子态密度导出：
在 $T=0$ 金属中，
以德布罗意波长 $\lambda$ 计数的电子分布。
\exsolution
含自旋简并的动量空间态密度满足
\[
 dn=\frac{4\pi p^2}{\pi^2\hbar^3}\frac{dp}{4}
 =\frac{p^2}{\pi^2\hbar^3}\,dp.
\]
利用 $p=h/\lambda=2\pi\hbar/\lambda$，
$|dp/d\lambda|=2\pi\hbar/\lambda^2$，得
\[
 \boxed{dn_\lambda=
 \frac{p^2}{\pi^2\hbar^3}
 \left|\frac{dp}{d\lambda}\right|d\lambda
 =\frac{8\pi}{\lambda^4}\,d\lambda},
 \quad \lambda\ge\lambda_F=h/p_F.
\]
\excheck 虽然分布不显含电子质量，
但是允许的最低波长
$\lambda_F=h/\sqrt{2m_eE_F}$包含材料电子浓度。

\moleculeproblem{204}{255}
零温金属钠的自由电子浓度
$n=2.5\times10^{22}\,\mathrm{cm^{-3}}$。
用理想气体动量通量表达式估算电子简并压强。
\exsolution
各向同性粒子气压强
$p=(1/3)n\langle pv\rangle$。
非相对论电子 $pv=2E$，故
\[
 p=\frac23n\overline E
 =\frac23n\frac35E_F=\frac25nE_F,
\]
\[
 \boxed{p=\frac{\hbar^2}{5m_e}(3\pi^2)^{2/3}n^{5/3}}.
\]
将 $n=2.5\times10^{28}\,\mathrm{m^{-3}}$ 代入，
得
$\boxed{p\simeq5.1\times10^9\,\mathrm{Pa}
 \approx5.0\times10^4\,\mathrm{atm}}$。
\excheck 零温量子简并压强不为零，
不同于经典 $p=nk_BT$ 的简单零温外推。

\moleculeproblem{205}{255}
电子管阴极温度从 $T=2000\,\mathrm K$
升高 $\Delta T=1.0\,\mathrm K$，
饱和发射电流增大 $\eta=1.4\%$。
求阴极材料逸出功。
\exsolution
由 Richardson--Dushman 关系
$I_s=C T^2e^{-W/(k_BT)}$，
对数微分
\[
 \frac{\Delta I_s}{I_s}\simeq
 \left(\frac2T+\frac W{k_BT^2}\right)\Delta T.
\]
故
\[
 \boxed{W=k_BT\left[
 \frac{\eta T}{\Delta T}-2\right]}.
\]
代 $\eta=0.014$ 得
$\boxed{W\simeq4.48\,\mathrm{eV}\approx4.5\,\mathrm{eV}}$。
\excheck $\eta$ 是相对增长的无量纲数 $0.014$，
而不是数字1.4。

\moleculeproblem{206}{255}
动能 $T=135\,\mathrm{eV}$ 的电子射入金属钠，
求其电子波的折射率。
设每钠原子贡献一个自由电子。
\exsolution
电子由真空进入金属时，
可用内势能降低 $U_0=E_F+W$，
其中 $E_F$ 是零温金属费米能，
$W$ 为金属表面逸出功。
动量分别
$p_{\rm out}=\sqrt{2m_eT}$、
$p_{\rm in}=\sqrt{2m_e(T+U_0)}$，故
\[
 \boxed{n_e=\frac{\lambda_{\rm out}}{\lambda_{\rm in}}
 =\sqrt{1+\frac{U_0}{T}}}.
\]
由 $n_{\rm Na}=\rho N_A/M_{\rm Na}$、
每原子一电子，
$E_F\simeq3.15\,\mathrm{eV}$。
采用钠逸出功 $W\simeq2.3\,\mathrm{eV}$
给 $U_0\simeq5.45\,\mathrm{eV}$，所以
$\boxed{n_e\simeq1.020\approx1.02}$。
\excheck 题面未另给逸出功，
故明确引用表值；数值近似
$n_e\simeq1+U_0/(2T)$ 而非 $1+U_0/T$。

\moleculeproblem{207}{255}
本征半导体从 $T_1=300\,\mathrm K$
升温到 $T_2=400\,\mathrm K$ 时，
电导率增为原来的 $\eta=5.0$ 倍。
忽略迁移率温度变化，
求形成电子--空穴对的最低能量。
\exsolution
本征载流子浓度的主指数为
$n_i\propto\exp[-E_g/(2k_BT)]$，
电导率同样满足
$\sigma\propto e^{-E_g/(2k_BT)}$。
于是
\[
 \ln\eta=\frac{E_g}{2k_B}
 \left(\frac1{T_1}-\frac1{T_2}\right),
\quad
 \boxed{E_g=\frac{2k_BT_1T_2}{T_2-T_1}\ln\eta}.
\]
代入
$\boxed{E_g\simeq0.333\,\mathrm{eV}}$。
\excheck 电子--空穴对的形成能与带隙 $E_g$
一致；电导率指数中有 $1/2$ 是因为本征
电子、空穴各贡献一半的激发熵关系。

\moleculeproblem{208}{255}
本征锗的低温光电截止波长为
$\lambda_{\rm th}=1.7\,\mu\mathrm m$。
估算室温时电阻率温度系数
$\alpha=(1/\rho)d\rho/dT$。
\exsolution
低温阈值对应带隙
$E_g=hc/\lambda_{\rm th}$。
忽略迁移率的温度依赖，
本征电阻率
$\rho(T)\propto e^{E_g/(2k_BT)}$。
因此
\[
 \boxed{\alpha=\frac{d\ln\rho}{dT}
 =-\frac{E_g}{2k_BT^2}
 =-\frac{hc}{2\lambda_{\rm th}k_BT^2}}.
\]
代 $T=300\,\mathrm K$，得
$\boxed{\alpha\simeq-0.047\,\mathrm{K^{-1}}
 \approx-0.05\,\mathrm{K^{-1}}}$。
\excheck 本征半导体升温载流子浓度增加，
电阻率下降，故温度系数为负。

\moleculeproblem{209}{256}
原图6.11以 $\ln\sigma$ 为纵轴，
$T^{-1}$ 为横轴，展示某 n 型半导体的
高温本征段与低温施主激活段。
根据两段直线斜率求禁带宽度和施主能级激活能。
\exsolution
本征高温段满足
$\ln\sigma=\mathrm{const}-E_g/(2k_BT)$，
因此曲线斜率
$s_i=d(\ln\sigma)/d(1/T)=-E_g/(2k_B)$。
低温施主热激活段
$\ln\sigma=\mathrm{const}-E_d/(2k_BT)$，
同理 $s_d=-E_d/(2k_B)$。
故
\[
 \boxed{E_g=-2k_Bs_i},\qquad
 \boxed{E_d=-2k_Bs_d}.
\]
由原图两段斜率约为
$s_i\simeq-7.0\times10^3\,\mathrm K$、
$s_d\simeq-3.5\times10^2\,\mathrm K$，
得到
$\boxed{E_g\simeq1.2\,\mathrm{eV},\quad
 E_d\simeq0.06\,\mathrm{eV}}$。
\excheck 原图横轴仅标 $T^{-1}$ 及1--4刻度，
实际读数需按常见 $10^{-3}\,\mathrm{K^{-1}}$
尺度理解；若按 $1\,\mathrm{K^{-1}}$ 读，
无法得到合理半导体能隙。源图隐含标度应在校订图中注明。

\moleculeproblem{210}{256}
本征半导体暗态电阻率
$\rho=50\,\Omega\cdot\mathrm{cm}$；
照光后降到
$\rho_1=40\,\Omega\cdot\mathrm{cm}$。
关光 $t=8\,\mathrm{ms}$ 后
变为 $\rho_2=45\,\Omega\cdot\mathrm{cm}$。
求非平衡电子与空穴平均寿命。
\exsolution
光生超额电导随关光时间按
$\Delta\sigma(t)=\Delta\sigma(0)e^{-t/\tau}$ 衰减。
$\sigma=1/\rho$，
故
\[
 e^{-t/\tau}=
 \frac{1/\rho_2-1/\rho}{1/\rho_1-1/\rho}
 =\frac{(\rho-\rho_2)\rho_1}
 {(\rho-\rho_1)\rho_2}.
\]
反解
\[
 \boxed{\tau=\frac t{
 \ln[(\rho-\rho_1)\rho_2/((\rho-\rho_2)\rho_1)]}
 \simeq9.87\,\mathrm{ms}\approx0.010\,\mathrm s}.
\]
\excheck 指数衰减的对象是光生\emph{超额电导}，
不是简单的电阻率差 $\rho(t)-\rho$。

\moleculeproblem{211}{256}
p型半导体板宽 $h=10\,\mathrm{mm}$、
长 $l=50\,\mathrm{mm}$，
磁感应强度 $B=5.0\,\mathrm{kG}$，
纵向电压 $V=10\,\mathrm V$、
霍尔电压 $V_H=50\,\mathrm{mV}$，
电阻率 $\rho=2.5\,\Omega\cdot\mathrm{cm}$。
求空穴数密度和迁移率。
\exsolution
纵向电场 $E_x=V/l$，
电流密度 $j=E_x/\rho=V/(\rho l)$。
霍尔电场 $E_H=V_H/h$。
单一空穴载流子的霍尔平衡
$E_H=jB/(ne)$，故
\[
 \boxed{n=\frac{hBV}{e\,l\rho V_H}}.
\]
迁移率由 $\sigma=ne\mu_h=1/\rho$ 给
\[
 \boxed{\mu_h=\frac1{ne\rho}
 =\frac{lV_H}{hBV}}.
\]
把 $B=0.50\,\mathrm T$、$\rho=0.025\,\Omega\mathrm m$
代入，
\[
 \boxed{n\simeq5.0\times10^{15}\,\mathrm{cm^{-3}},
 \quad\mu_h\simeq0.050\,\mathrm{m^2/(V\,s)}}.
\]
\excheck 电阻率 $2.5\,\Omega\cdot\mathrm{cm}$
必须换成 $0.025\,\Omega\mathrm m$。

\moleculeproblem{212}{256}
在本征锗中，外磁场
$B=5.0\,\mathrm{kG}$；
霍尔横向电场强度比纵向电场
低 $\eta=10$ 倍。求电子、空穴迁移率之差。
\exsolution
本征半导体电子与空穴浓度相等。
低磁场双载流子模型
\[
 R_H=\frac{\mu_h^2-\mu_e^2}
 {en(\mu_h+\mu_e)^2}
 =\frac{\mu_h-\mu_e}{en(\mu_h+\mu_e)}.
\]
电导率
$\sigma=en(\mu_h+\mu_e)$，因此
\[
 \frac{E_H}{E_x}=R_H\sigma B
 =(\mu_h-\mu_e)B.
\]
取绝对值及 $E_H/E_x=1/\eta$，
\[
 \boxed{|\mu_h-\mu_e|=\frac1{\eta B}
 =\frac1{10\times0.50}
 =0.20\,\mathrm{m^2/(V\,s)}}.
\]
\excheck 在本征样品中
不能套用单载流子的 $R_H=1/(ne)$，
必须使用双载流子霍尔表达式。

\moleculeproblem{213}{256}
某半导体未观察到霍尔效应，
其电子迁移率是空穴迁移率的
$\eta=2.0$ 倍。
求空穴数密度与导带电子数密度之比。
\exsolution
双载流子霍尔系数
\[
 R_H=\frac{n_h\mu_h^2-n_e\mu_e^2}
 {e(n_h\mu_h+n_e\mu_e)^2}.
\]
若霍尔效应抵消，$R_H=0$，
故
$n_h\mu_h^2=n_e\mu_e^2$，
从而
\[
 \boxed{\frac{n_h}{n_e}=
 \left(\frac{\mu_e}{\mu_h}\right)^2
 =\eta^2=4}.
\]
\excheck 霍尔效应为零不说明无载流子，
而是电子和空穴的霍尔贡献恰好相互抵消。
